% 11. kapitulua - Ohiko arazoak eta errezetak
% Helburua: aurreko liburuko 10. kapituluaren formatua (arazoa -> errezeta), agenteekin lan egitean sortzen direnetarako:
%           agenteak "eginda" dio baina testak huts; agenteak testak ahuldu ditu; agenteen adarren arteko gatazkak;
%           testuingurua galdu; kostua lehertu...
% Aurreko liburutik berrerabil daitekeena:
%   - aurreko_liburua_2024/kapituluak/10-arazoak.tex (egitura eta tonua)
% Irudiak: Pictures/11/ (10.arazoen_grafoa, mind_blower; goiburu hutsa header_11.pdf.xoj)
\todo[inline]{Laburpena (behin-behinekoa): Aurreko eskuliburuko azken kapituluaren formatu berean, agenteekin lan egitean sortzen diren arazo ohikoenak bilduko ditugu, bakoitza errezeta batekin: agenteak ``eginda'' dio baina testek huts egiten dute; agenteak testak aldatu ditu pasa zitezen; bi agentek fitxategi bera aldatu dute eta gatazka sortu da; agenteak zereginaren mugak gainditu ditu eta behar ez zena ukitu du; zeregina lausoegia zen eta agenteak asmatu egin du; berrikuspenen ilara pilatu da eta fluxua gelditu da; agenteak testuingurua galdu du zeregin luze batean; kostua espero baino askoz handiagoa izan da; agenteak sekretu bat edo baimen gehiegi erabili ditu. Errezeta bakoitzak arazoa antzemateko seinaleak, konponbidea eta prebentzioa ditu.}
